\bookchapter{Capítulo 19: Solução de problemas e autoajuda}

\chapterintro{O celular, assim como as pessoas, de vez em quando também “adoece”. Travamentos, fechamentos inesperados, tela preta, não carrega... A maioria desses problemas não é tão assustadora quanto parece. Você pode perfeitamente tentar resolvê-los sozinho antes de levar ao conserto. Este capítulo é o seu “manual de pronto-socorro para celular” — primeiro tente os métodos de baixo custo; se não funcionar, procure um técnico profissional.}

\booksection{I. Celular lento e aquecendo — a pergunta mais frequente}

Sinceramente, é supernormal o celular começar a ficar lento depois de cerca de um ano de uso. A causa nem sempre é defeito; na maioria das vezes é “acúmulo de sujeira” — arquivos inúteis, processos em segundo plano, espaço de armazenamento gradualmente preenchido.

\booksubsection{Primeiro faça o diagnóstico, não se apresse em consertar:}

\begin{enumerate}
  \item \textbf{Verifique o espaço de armazenamento}: Vá em Ajustes > Armazenamento e veja o espaço livre. Se estiver abaixo de 10\%, o sistema fica extremamente lento, porque está lutando para liberar espaço. Minha experiência: mantenha pelo menos 15\%-20\% de espaço livre para o celular funcionar com fluidez.
  \item \textbf{Verifique o ranking de consumo de bateria}: Nas configurações de bateria, veja o ranking de consumo das últimas 24 horas. Se algum app está consumindo muito, mas você quase não o usa, é ele rodando escondido em segundo plano. Desinstale ou restrinja a atividade em segundo plano.
  \item \textbf{Pense no que você fez recentemente}: O celular ficou lento depois de atualizar o sistema? (É normal ficar lento nos primeiros dias após uma grande atualização de versão, pois o sistema está reindexando em segundo plano). Ou ficou lento depois de instalar um novo app? (Desinstale e teste).
  \item \textbf{Observe o ambiente de uso}: Se o celular ficou exposto ao sol forte, com capa grossa e carregando enquanto joga, o superaquecimento é inevitável. O calor faz o CPU reduzir automaticamente a frequência, e você sente o celular de repente mais lento. Deixe esfriar e teste novamente.
\end{enumerate}

\textbf{Métodos de tratamento (do simples ao complexo)}: Reinicie o celular (resolve 60\% dos problemas) Libere espaço de armazenamento (exclua apps e vídeos inúteis) Se um app específico está causando o problema, limpe o cache dele ou desinstale e reinstale Último recurso: faça backup dos dados e restaure as configurações de fábrica.

\booksection{II. App fecha sozinho, não responde — o app não está obedecendo}

Um app que fecha sozinho ao abrir ou trava de repente durante o uso. O fluxo de tratamento é como uma consulta médica, passo a passo:

\begin{enumerate}
  \item \textbf{Primeiro feche e reabra}: Remova o app dos apps recentes (deslize para cima) e abra novamente. A maioria dos bugs temporários se resolve assim.
  \item \textbf{Atualize o app}: Vá na loja de aplicativos e veja se há atualização. Às vezes é porque a versão do app está muito antiga e incompatível com o sistema.
  \item \textbf{Verifique as permissões}: Alguns apps precisam de certas permissões para funcionar corretamente, mas você não as concedeu. Por exemplo, um app de câmera sem permissão de câmera fecha ao abrir. Vá nas configurações, encontre o app e conceda as permissões necessárias.
  \item \textbf{Limpe o cache}: Vá em Ajustes > Gerenciamento de apps > Encontre o app > Armazenamento > Limpar cache. Atenção: é "Limpar cache", não "Limpar dados"! Limpar dados apaga o estado de login e as configurações do app, equivalente a restaurar o app ao estado de fábrica.
  \item \textbf{Desinstalar e reinstalar}: Último recurso. Antes de desinstalar, confirme se os dados do app serão perdidos — no caso do WeChat, faça backup do histórico de conversas primeiro; no caso de jogos, confirme se a conta está vinculada a um número de telefone.
\end{enumerate}

\booksection{III. Não carrega ou carrega devagar — pode não ser culpa da bateria}

Não conseguir carregar é um problema que gera muita ansiedade, mas muitas vezes a causa está nos acessórios, não no celular.

\begin{itemize}
  \item \textbf{Primeiro verifique a porta de carregamento}: Com um palito de dente ou uma ferramenta não condutiva, limpe delicadamente a poeira e os fiapos da porta de carregamento. Eu mesmo já passei por isso — um fio de linha do bolso entupiu exatamente a porta de carregamento. Ao conectar, parecia mau contato. Tirei o fiapo e resolveu.
  \item \textbf{Teste com outro cabo e carregador}: A taxa de danos dos cabos é muito maior que a dos carregadores, especialmente cabos que são enrolados com frequência. Teste com um cabo e carregador que você sabe que funcionam, para identificar rapidamente qual acessório está com defeito.
  \item \textbf{Problemas com carregamento sem fio}: Se você usa carregamento sem fio, verifique se a capa do celular é muito grossa ou se há metal atrás (como suporte magnético). Além disso, posicione o celular exatamente no centro da base de carregamento — se estiver torto, pode não carregar.
  \item \textbf{Carregamento lento}: Confirme se a potência do carregador é suficiente (iPhone precisa de pelo menos 20W com carregador PD para carregamento rápido), se o cabo suporta carregamento rápido (nem todo cabo aguenta corrente acima de 3A) e se o celular está muito quente (muito calor reduz automaticamente a potência para proteger a bateria).
\end{itemize}

\booksection{IV. Rede não conecta, Wi-Fi vive caindo}

Na rua, os dados móveis funcionam bem. Em casa, o Wi-Fi não conecta. Ou o Wi-Fi está conectado, mas não tem internet. Siga esta ordem de verificação:

\begin{enumerate}
  \item \textbf{Ative e desative o Modo Avião}: Deslize para baixo no Centro de Controle, ative o Modo Avião, espere 5 segundos e desative. Isso equivale a fazer o celular procurar rede novamente. Muito eficaz.
  \item \textbf{"Esqueça" a rede Wi-Fi e reconecte}: Nas configurações de Wi-Fi, toque no ícone de exclamação ao lado do nome da rede (Android é uma engrenagem), selecione "Ignorar esta rede" ou "Esquecer" e depois digite a senha novamente para conectar.
  \item \textbf{Reinicie o roteador}: Muita gente não desliga o roteador de casa há meses. Com o tempo, ele fica instável. Desconecte da tomada, espere 30 segundos e conecte novamente. A maioria dos problemas de Wi-Fi se resolve assim.
  \item \textbf{Verifique se é problema do celular}: Outros dispositivos conectam ao mesmo Wi-Fi normalmente? Se todos funcionam bem e só o seu celular não, então o celular pode precisar ser reiniciado ou ter as configurações de rede redefinidas (Ajustes > Geral > Transferir ou Redefinir > Redefinir Ajustes de Rede — isso não apaga seus dados, mas limpa todas as senhas de Wi-Fi salvas).
\end{enumerate}

\booksection{V. Problemas de tela: toque não responde, tela piscando, marcas na tela}

\begin{itemize}
  \item \textbf{Toque pouco responsivo}: Primeiro, remova a película de vidro temperado para ver se é culpa dela — películas baratas, com o tempo, realmente afetam o toque. Água na tela ou gordura nos dedos também causam desvio no toque. Se nenhuma dessas for a causa, reinicie e teste novamente.
  \item \textbf{Marcas/imagens fantasmas na tela (burn-in)}: Exibir a mesma imagem por muito tempo (como assistir TikTok com a tela sempre ligada) pode deixar "imagens fantasmas" na tela OLED. Isso não significa que a tela está quebrada, é apenas temporário. Desligue a tela e deixe descansar por meia hora; geralmente se recupera.
  \item \textbf{Tela piscando}: Se piscar em baixo brilho, pode ser problema de modulação PWM, comum em telas OLED. Procure nas configurações por "DC Dimming" ou "Anti-flicker" e ative — melhora muito. Nem todos os celulares têm essa função.
  \item \textbf{Uma área da tela completamente preta ou clara}: Isso é mais grave. Provavelmente é dano físico na tela e precisa ser trocada. Telas trocadas em assistências não oficiais geralmente têm cores e brilho inferiores ao original e comprometem a resistência à água.
\end{itemize}

\booksection{VI. Não liga, trava na tela do logo — primeiro, não entre em pânico}

Esta situação é a que mais acelera o coração, mas mantenha a calma. Primeiro tente isto:

\begin{itemize}
  \item \textbf{Carregue por meia hora antes de qualquer coisa}: Muitas vezes é apenas a bateria completamente esgotada. Conecte o carregador original e espere pelo menos 30 minutos antes de tentar ligar. Alguns celulares, depois de descarregar totalmente, precisam atingir um certo nível de carga para reagir.
  \item \textbf{Reinicialização forçada}: iPhone com tela cheia: pressione rapidamente Aumentar Volume, depois Diminuir Volume, e mantenha pressionado o botão lateral (ligar/desligar) até aparecer o logo da Apple (pode levar de 10 a 15 segundos). Método universal Android: mantenha pressionado o botão liga/desliga por mais de 10 segundos, até vibrar ou a tela acender.
  \item \textbf{Travado na tela de inicialização, reiniciando repetidamente}: Esta situação é chamada de "maçã branca" (iPhone) ou "trava no logo" (Android). Se a reinicialização forçada não funcionar, pode ser necessário entrar no modo de recuperação e reparar pelo computador. iPhone: conecte ao computador, abra iTunes/Finder, entre no modo de recuperação e selecione "Atualizar" (não "Restaurar") — reinstala o sistema mantendo os dados. Android: entre no modo Fastboot ou Recovery. As combinações de teclas variam por marca; recomenda-se pesquisar "modo de recuperação + modelo do seu celular".
  \item \textbf{Se nada acima funcionar}: Só resta procurar um reparo profissional. Se os dados estiverem com backup, não precisa se preocupar tanto.
\end{itemize}

\booksection{VII. Problemas de som: alto-falante não funciona, volume baixo nas chamadas}

\begin{itemize}
  \item \textbf{Alto-falante sem som}: Primeiro, verifique se não está no modo silencioso (chave de mudo na lateral do iPhone), no modo Não Perturbe ou com o volume no mínimo. Parece bobo, mas isso representa 40\% dos casos. Depois, verifique se o Bluetooth está conectado a algum fone ou caixa de som.
  \item \textbf{Durante a chamada, a outra pessoa não ouve você}: Verifique se o microfone não está obstruído — os furinhos na parte inferior do celular são facilmente entupidos por poeira. Limpe delicadamente com uma escova macia. Não use agulha ou palito de dente para cutucar com força, pois pode danificar a tela de proteção contra poeira.
  \item \textbf{Som rouco/distorcido}: Se for causado por entrada de água no alto-falante — não use secador de cabelo com ar quente! Seque a superfície com papel toalha e deixe secar naturalmente. Dizem na internet que apps de "expulsão de água por ondas sonoras" têm certo efeito (emitem sons de baixa frequência para vibrar a água para fora). Pode tentar.
\end{itemize}

\booksection{VIII. Procurar assistência oficial e evitar as armadilhas}

Se todos os métodos acima foram tentados e não funcionaram, é hora de procurar um reparo profissional. Mas o mercado de reparos é cheio de armadilhas. Vou contar as que eu mesmo já enfrentei e o que aprendi:

\begin{itemize}
  \item \textbf{Procure a assistência oficial de verdade, não a que "parece oficial"}: Vá ao site oficial da marca (Apple.com, Huawei.com etc.) e encontre a lista de "Prestadores de Serviço Autorizados". Lojas não autorizadas com um logo da Apple na fachada podem não usar peças originais. Depois do reparo, a resistência à água do seu celular já era.
  \item \textbf{O que fazer antes de levar ao reparo}: Faça backup de todos os dados (não conte que os dados não serão perdidos no reparo), saia de todas as contas importantes (Apple ID, WeChat, Alipay), desative o "Buscar meu celular". Alguns pontos de reparo pedem a senha de bloqueio de tela — isso tem risco de privacidade. É melhor você mesmo sair das contas antes de entregar.
  \item \textbf{Armadilhas na troca de tela}: Telas não oficiais se dividem em "tela original com vidro trocado", "tela com vidro paralelo" e "tela montada" (toda paralela). Os preços variam muito, e a qualidade de exibição e durabilidade também. Se o celular ainda estiver na garantia, vá direto para a assistência oficial.
  \item \textbf{Armadilhas na troca de bateria}: Baterias não originais podem ter capacidade falsa e, depois de meio ano, degradam mais do que uma bateria original degradaria em dois anos. O preço da troca de bateria na Apple pela assistência oficial não é muito mais caro que o paralelo. Recomenda-se ir na oficial.
  \item \textbf{Armadilhas na recuperação de dados}: Depois que o celular sofre danos por água/queda, se a assistência disser que "recuperação de dados tem custo extra", cuidado — alguns estabelecimentos ruins aumentam propositalmente o grau do dano para cobrar mais. Recuperação de dados profissional é: primeiro o orçamento, depois o serviço. E se não tiver sucesso, não cobra.
\end{itemize}

\begin{tcbwarn}{A primeira dica universal: reinicie o celular}
Não importa qual problema você encontre — travamento, fechamento inesperado, rede que não conecta, som estranho — primeiro reinicie o celular. Esta frase parece enganação, mas, pela minha experiência, pelo menos 60\% dos pequenos problemas se resolvem depois de reiniciar. O celular, assim como o computador, com o tempo de uso acumula várias coisas bagunçadas na memória. Reiniciar é limpar tudo e recomeçar. Então, da próxima vez que o celular der problema, não se apresse em pesquisar “o que fazer” — reinicie primeiro e depois veja.

\end{tcbwarn}

\begin{tcbpractice}{Anote agora: estas operações de emergência devem estar na ponta da língua}
\begin{enumerate}
  \item \textbf{Reinicialização forçada}: iPhone com tela cheia: Aumentar Volume Diminuir Volume segurar botão liga/desliga até aparecer o logo. Android: segurar botão liga/desliga por 10-15 segundos até vibrar.
  \item \textbf{Limpar porta de carregamento}: Pegue uma ferramenta pequena (palito de dente ou chave de ejeção do chip) e limpe delicadamente a poeira da porta de carregamento.
  \item \textbf{Reconectar Wi-Fi}: Nas configurações de Wi-Fi, "esqueça" sua rede atual e digite a senha novamente para conectar.
  \item \textbf{Ativar/desativar Modo Avião}: Deslize para baixo no Centro de Controle, ative o Modo Avião, conte mentalmente 5 segundos, desative e veja se a rede voltou.
  \item Tire print dos passos dessas operações ou anote nas Notas. Quando encontrar um problema, não precisa pesquisar desesperadamente — é só abrir as Notas.
\end{enumerate}

\end{tcbpractice}
